{\small
\setlength{\tabcolsep}{5pt}
\begin{longtable}{@{}>{\raggedright\arraybackslash}p{0.48\linewidth}>{\raggedright\arraybackslash}p{0.48\linewidth}@{}}
\toprule
% DOCX table after paragraph 1168, row 0, column 0
\textbf{EM Waves}
&
% DOCX table after paragraph 1168, row 0, column 1
\textbf{Gravitational Waves}
\\[6pt]
\midrule
% DOCX table after paragraph 1168, row 1, column 0
They are oscillations of electric and magnetic fields that propagate through spacetime.
&
% DOCX table after paragraph 1168, row 1, column 1
They are oscillations of the spacetime itself.
\\[6pt]
% DOCX table after paragraph 1168, row 2, column 0
The simplest are those relative to dipole\index{dipole!radiation} radiation, i.e., generated by an oscillating dipole\index{dipole!electric} (two opposite electric charges).
&
% DOCX table after paragraph 1168, row 2, column 1
The simplest are those relative to quadrupolar radiation\index{quadrupole!radiation}, i.e., generated by an oscillating mass quadrupole\index{quadrupole!mass} (a non-circular distribution of mass).
\\[6pt]
% DOCX table after paragraph 1168, row 3, column 0
They are almost always the result of the incoherent superposition of contributions from millions of electrons, atoms or molecules.
&
% DOCX table after paragraph 1168, row 3, column 1
They are produced by the coherent movement of large amounts of mass and energy, be they in the form of astronomical objects or energy condensations.
\\[6pt]
% DOCX table after paragraph 1168, row 4, column 0
They show us the details of the thermodynamic state\index{thermodynamic state} (temperature, density) of the sources that produce them.
&
% DOCX table after paragraph 1168, row 4, column 1
They show us the details of the dynamics and global properties (size, mass, and velocity) of large concentrations of mass/energy.
\\[6pt]
% DOCX table after paragraph 1168, row 5, column 0
Their wavelengths are very short compared to the size of the astronomical sources. Hence, it is possible to obtain ``images'' of the sources.
&
% DOCX table after paragraph 1168, row 5, column 1
Their wavelengths are often comparable to the size of the sources that produce them, if not greater. Hence, it is not possible to produce ``images'' of the sources.
\\[6pt]
% DOCX table after paragraph 1168, row 6, column 0
They suffer from absorption, diffusion and dispersion by the matter they encounter on their path.
&
% DOCX table after paragraph 1168, row 6, column 1
They propagate undisturbed, at the speed of light, through all types of materials.
\\[6pt]
% DOCX table after paragraph 1168, row 7, column 0
Those emitted by celestial objects have a spectrum\index{spectrum!electromagnetic} that is essentially independent of the mass of the source; for instance, the emissions from stellar and supermassive black holes are very similar.
&
% DOCX table after paragraph 1168, row 7, column 1
Those emitted by celestial objects have a spectrum\index{spectrum!gravitational wave} that depends on the mass of the source: the frequency of gravitational waves from stellar black holes is much higher than that from supermassive black holes.
\\[6pt]
% DOCX table after paragraph 1168, row 8, column 0
They have wavelengths shorter than the curvature radii of the objects they encounter as they propagate. For this reason, they experience gravitational lensing\index{gravitational!lensing}, i.e., distortion and amplification.
&
% DOCX table after paragraph 1168, row 8, column 1
They have wavelengths comparable if not greater than the curvature radii of the objects they encounter while propagating. For this reason, gravitational lensing is less pronounced and more difficult to measure.
\\[6pt]
\bottomrule
\end{longtable}
}
